% 本文件由 LaTeX 排版 Agent 根据有效 translation.json 自动生成。
% author=Hippolytus; work_id=0521; title=驳诺厄特

\chapter{驳诺厄特}

% structure: p0001 source=h1 target=omitted reason=page-title-already-rendered-by-parent

1. 另有一些人暗中传授另一种道理，他们成了一名诺厄特的弟子，此人原籍\Place{斯米纳}，生活于不久以前。他极度骄傲自大，受到一种怪异精神的虚妄所驱使。他声称基督就是圣父本身，圣父本身受生、受难、死亡。你看，骄傲之心和一种多么怪异的膨胀精神已潜入他心中。从他所行的其他事，我们已得到证明：他说话不是出于纯洁的精神；因为亵渎圣神的人被排除在圣继承之外。他自称是梅瑟，而他的兄弟是亚郎。当有福的诸位长老听闻此事，便将他召至教会面前，审问他。但他起初否认持有这些意见。然而后来，他躲到一些人中，并聚集了一些接受同一谬误的人，此后便希望公开坚持自己的教义为正确。有福的长老们再次将他召至面前审问。但他反对他们说：\Quote{那么，我光荣基督，做了什么恶事？} 长老们回答他：\Quote{我们确实也只认识一位天主；我们认识基督；我们知道圣子受苦正如祂受苦，死亡正如祂死亡，第三日复活，坐在圣父右边，并来审判生者死者。我们所说的这些是我们所学的。} 于是审问他之后，他们将他逐出教会。他竟骄傲到如此地步，以致开办了一所学校。\par

2. 如今他们寻求通过引用法律中的话为自己的教义提供根据：\Quote{我是你们祖先的天主：你们在我面前不可有别的神}；又在另一处：\Quote{我是元始，}祂说，\Quote{也是终结；除我以外没有别的神。} 因此他们说他们证明天主是独一的。然后他们这样回答：\Quote{因此，如果我承认基督是天主，那么祂就是圣父本身，如果祂确实是天主；基督受苦，祂自己就是天主；因此圣父受苦，因为祂就是圣父本身。} 然而情况并非如此；因为圣经并非这样陈述此事。但他们还引用其他证言，说：经上写道：\Quote{这就是我们的天主，没有别的可比拟祂。祂寻得知识的一切道路，赐给祂的仆人雅各伯和祂所爱的以色列。此后祂在地上显现，与人同住。} 那么你看，他说，这就是那位独一的天主，祂后来显现并与人同住。又在另一处他说：\Quote{埃及劳碌；埃塞俄比亚和色巴的魁梧人物要来到你这里（他们必作你们的奴仆），他们将在你后面带着镣铐而来，他们向你下拜，因为天主在你中间；他们向你祈求：除你以外没有别的神。因为你是天主，我们未曾认识；以色列的天主，救世主。} 你看，他说，圣经如何宣告独一的天主？既然这点清楚显明，这些经文是为此作证，他说，我不得不承认独一的天主，并让这独一者受苦。因为基督是天主，为我们受苦，祂自己就是圣父，为能拯救我们。我们不能有别的说法，他说；因为宗徒也承认独一天主，当他说：\Quote{圣祖们也是他们的，按血统基督也是从他们而来的，祂是在万有之上，永远可赞美的天主。}\par

3. 因此，他们选择这样陈述这些事，并且只使用一类经文；就像特奥多图斯曾用来证明基督是纯粹的人所用的片面方式一样。但双方都没有正确理解这事，因为圣经本身就驳斥了他们的愚昧，并为真理作证。弟兄们，看他们引入了多么鲁莽和大胆的教义，当他们无耻地说：圣父就是基督，就是圣子，祂自己受生，自己受苦，自己使自己复活。但并非如此。圣经讲的是正确的；但诺厄特与他们不同心。然而，尽管诺厄特不理解真理，圣经不应因此被否定。因为谁会不说只有一位天主？但他不会因此否认\OriginalTerm{安排}{economy}（即天主圣三的位格数目和次序）。因此，处理这个问题的正确方式是首先驳斥这些人对这些经文所做的解释，然后说明它们的真正含义。因为正确的做法是首先阐明真理：圣父是独一天主，\Quote{一切家族都出于祂}，\Quote{万有都借着祂，万有都归于祂，我们也在祂内}。\par

4. 如我所说，我们来看看他是如何被驳倒的，然后就陈述真理。现在他引用了这些话：\BibleQuote{埃及劳碌，古实和赛巴人的货物}等等，直到\BibleQuote{因为你是以色列的天主，拯救者}。而他引用这些话时并不明白它们的上文。因为每当他们图谋暗中行事时，就割裂圣经。但让他引用整段经文吧，他就会发现写作时所保持的理由。因为我们在稍向上之处有这一小段的开始；我们当然应该从那里开始，指明这话是对谁说的、关于谁的。因为上面，这一小段的开始是这样写的：\BibleQuote{你们可以问我关于我的儿子和我的女儿，并关于我手的工作，命令我。我造了大地，并造了人在地上；我亲手铺张了诸天；我命令了所有星辰。我凭公义兴起了他，他的一切道路都是正直的。他必建造我的城，并要释放被掳的人；不为代价，也不为报酬，这是万军的上主说的。万军的上主这样说：埃及劳碌，古实和赛巴人的货物，身材高大的人，必都投降你，也要作你的奴仆；他们必带着锁链过来，并要向你下拜；他们要恳求你，因为天主在你中间；除你以外没有别的神。因为你是天主，我们不知道；以色列的天主，拯救者}。\Quote{因此，在你之内，}他说，\Quote{有天主。}但除了在基督耶稣——圣父的圣言和\OriginalTerm{救恩计划的奥迹}{mystery of the economy}——之内，天主在哪里呢？再次，为显明关于他的真理，他指出了他在肉身中的事实，当他说：\BibleQuote{我凭公义兴起他，他的一切道路都是正直的。}因为这是什么？圣父如此为谁作证？正是为圣子，圣父说：\BibleQuote{我凭公义兴起他。}而圣父确实凭公义兴起了他的圣子，宗徒保禄作证说：\BibleQuote{但那使基督耶稣从死人中复活者的圣神若住在你们内，那使基督从死人中复活的，也必借着他住在你们内的圣神，使你们必死的身体活过来。}看，先知的话就这样成就了：\BibleQuote{我凭公义兴起他。}而他说\BibleQuote{天主在你中间}，是指\OriginalTerm{救恩计划的奥迹}{mystery of the economy}，因为当圣言降生成人、成为人时，圣父在圣子内，圣子也在圣父内，同时圣子生活在人中间。因此，弟兄们，这正是所预表的：实际上，借着圣神和童贞玛利亚，\OriginalTerm{救恩计划的奥迹}{mystery of the economy}就是这位圣言，构成天主唯一的圣子。这不是我随意说的，而是那自天降下的亲自见证；因为他这样说：\BibleQuote{没有人升过天，除了那从天降下的、人子，他是在天上的。}那么他还能寻求这所写的之外什么呢？他难道要说，肉身上了天吗？然而有那由圣父的圣言作为祭品呈现的肉身——那由圣神和童贞玛利亚而来的肉身，（并被）显明为天主的完美圣子。因此，显然他将自己献给了圣父。在此以前，天上没有肉身。那么，谁在天上呢？岂不是那\OriginalTerm{未降生成人的圣言}{Word unincarnate}，他被派遣以表明他在地上同时也在天上？因为他是圣言，他是圣神，他是能力。他取了那在人间常见常用的名字，并从起初因着他将要成为的被称为人子，尽管他尚未成为人，正如达尼尔作证说：\BibleQuote{我观看，见有一位像人子的，乘着天上的云彩而来。}因此，他正确地说，那在天上的从起初就以此名被称呼，即天主的圣言，因为他从起初就是那一位。\par

5. 但在另一段经文中又是什么意思呢，他说：\BibleQuote{这是天主，再没有别的能与他相比？}他说得对。因为谁能与圣父相比呢？但他又说：\BibleQuote{这是我们的天主，再没有别的能与他相比。他找出了全部知识之路，并把它赐给了他的仆人\Person{雅各伯}和他所爱的\Person{以色列}。}他说得好。因为谁是\Person{雅各伯}他的仆人，他所爱的\Person{以色列}，岂不是他呼喊的那一位：\BibleQuote{这是我的爱子，我所喜悦的，你们要听从他？}然后，从圣父接受了全部知识的那\OriginalTerm{完美的以色列}{perfect Israel}、真\Person{雅各伯}，后来在地上显现，并与人交往。而以色列又指谁呢，岂不就是\Emph{看见天主的人}？除了圣子之外，没有人看见天主；他是完美的人，唯有他宣告圣父的旨意。因为若望也说：\BibleQuote{从来没有人见过天主；只有在圣父怀里的独生子，他彰显了天主。}又说：\BibleQuote{那从天降下的，见证他所听见和看见的。}这就是那圣父将全部知识赐给他的人，他在地上显现，与人交往。\par

6. 让我们接着看宗徒的话：\BibleQuote{列祖是他们的，按肉身说，基督也是从他们来的；祂是在万有之上，永远受赞美的天主。} 这话正确地、清楚地宣告了真理的奥迹。那在万有之上的就是天主；因为祂大胆地说：\BibleQuote{我父将一切都交给了我。} 那在万有之上、永远受赞美的天主，已经降生成人；祂成为人，却〔仍然是〕永远的天主。为此若望也说：\BibleQuote{那今在、昔在、将来永在的全能者。} 祂恰当地称基督为全能者。因为在这话里，祂所说的正是基督自己所见证的。因为基督作证说：\BibleQuote{我父将一切都交给了我。} 基督掌管万有，并由圣父被立为全能者。同样，保禄在阐明万有都交给了祂这一真理时，也说：\BibleQuote{基督是初果；以后，在基督再来时，是那些属于基督的人。然后是终结，那时祂将把王国交给天主父；那时祂将消灭一切统治者、一切掌权者和有能者。因为祂必须为王，直到把一切仇敌放在祂脚下。最后被毁灭的仇敌是死亡。因为万有都被放在祂脚下。但祂说\InnerQuote{万有都放在祂脚下}时，明显那将万有放在祂脚下的不在内。那时，圣子也要顺服那将万有放在祂脚下的，好使天主成为万有中的万有。} 因此，如果万有都放在祂脚下，只有那放万有在祂脚下的除外，那么祂是万有的主，圣父是祂的主，好使在一切中彰显独一天主，万有连同基督都服从于祂，圣父已将万有服从于基督，除了祂自己。而这正是基督亲自说的，如同在福音中祂承认祂是自己的父和自己的天主。因为祂如此说：\BibleQuote{我升到我的父和你们的父那里，到我的天主和你们的天主那里。} 如果诺厄托胆敢说祂就是圣父自己，那么根据福音的话，他还能说基督升到哪个父那里呢？但如果他希望我们放弃福音而相信他的胡言乱语，那他就白费力气了；因为\BibleQuote{我们必须服从天主而不是服从人。}\par

7. 如果他再引用祂自己的话：\BibleQuote{我与父原是一体。} 让他注意这个事实，并明白祂并没有说\BibleQuote{我与父是\Emph{一（单数）}}，而是说\BibleQuote{（我们）是\Emph{一（复数）}}。因为\Quote{是（复数）}这个词不是指一个位格，而是指两个位格和一个权能。祂亲自使这一点清楚，当祂向圣父谈及门徒时说：\BibleQuote{你赐给我的荣耀，我已给了他们，为使他们合而为一，如同我们原为一体：我在他们内，你在我内，使他们完全合而为一，为使世界知道是你派遣了我。} 诺厄托派对此有何话可说？难道在本质上，所有都是一体吗？还是我们在权能和合一心思的状态中成为一体？同样，圣子被派遣，却不被世上的人所认识，祂宣认自己以权能和心意在父内。因为圣子是圣父的唯一心意。我们拥有圣父心意的人（这样）相信祂；但那不拥有的人否认了圣子。如果，他们再举出斐理伯询问父的事，说：\BibleQuote{请把父显示给我们，我们就满足了。} 主这样回答他：\BibleQuote{斐理伯，这么长久的时间，我和你们在一起，你还不认识我吗？看见我的，就是看见了父。你不信我在父内，父在我内吗？} 如果他们认为这段经文确认了他们的教义，好像祂承认自己就是父，让他们知道这经文恰好反对他们，他们正被这句话驳倒。因为虽然基督论及自己，并在众人面前显明自己是圣子，他们仍未认识祂是这样，也不能领会或默观祂的真实权能。斐理伯未能接受这一点（就当时可见的程度而言），请求看见父。主于是对他说：\BibleQuote{斐理伯，这么长久的时间，我和你们在一起，你还不认识我吗？看见我的，就是看见了父。} 祂的意思是说：如果你看见了我，你就可以透过我认识父。因为透过相似（原像）的肖像，父得以被容易认识。但如果你不认识肖像——就是圣子，你怎能寻求看见父呢？这一点由本章其余部分证明，那表明圣子\BibleQuote{受派遣，来自父，又回到父那里。}\par

8. 许多其他经文，更确切地说，所有经文都见证真理。因此，即使人不愿意，他也被迫承认全能的天主圣父、及天主子耶稣基督——祂是天主，降生成人，圣父也把一切置于祂之下，只有自己除外——并承认圣神；所以，这些是三位。但如果他渴望了解如何仍然表明只有一位天主，让他知道祂的权能是一个。因此，就权能而言，天主是独一的。但就\OriginalTerm{上智的安排}{\GreekText{οἰκονομία}}而言，有三重的显示，正如以后我们阐述真道时将要证明的。然而，在这些由我们所阐述的事上，我们是一致的。因为有一位天主，我们必须相信祂，祂是无始、无受苦难、不死不灭的，按祂的意愿、按祂的方式、按祂的时间行一切事。那么，这个对真理一无所知的诺厄托，还敢对这些事说什么呢？既然诺厄托已被驳倒，现在让我们转向真理本身的展示，好使我们确立真理，反对这一切兴起的强大异端，而它们无法提出任何中肯的论点。\par

9. 弟兄们，只有一位天主，我们对祂的认识来自圣经，而非其他来源。正如一个人，若想通晓今世的智慧，除了掌握哲学家的学说外，别无他法；同样，我们所有愿意实践虔诚的人，除了天主的圣言外，也无法从其他途径学习其实践。因此，凡圣经所宣告的，让我们注视；凡圣经所教导的，让我们学习；如圣父愿意我们相信的，让我们相信；如祂愿意圣子受荣耀，让我们荣耀祂；如祂愿意圣神被赐予，让我们领受祂。不是按照我们自己的意愿，也不是按照我们自己的心意，更不是强行使用天主所赐予的事物，而是按照祂选择通过圣经教导我们的方式，让我们如此辨别。\par

10. 天主独自存在，没有与祂同时代的事物，决定创造世界。祂在心中构思世界，愿意并说出圣言，就创造了它；它立刻出现，按祂所喜悦的样式形成。因此，对我们来说，知道没有任何与天主同时代的事物就足够了。除了祂，没有别的；然而祂独自存在，却存在于多样性中。因为祂并非没有理性、智慧、能力或谋略。万有都在祂内，祂就是万有。当祂愿意时，并如祂所愿，祂在祂所定的时间内彰显了祂的圣言，并借着祂创造了万有。当祂愿意时，祂就行动；祂思想时，祂就执行；祂说话时，祂就彰显；祂创造时，祂以智慧设计。因为祂以理性和智慧形成一切受造之物——在理性中创造它们，在智慧中安排它们。于是，祂按祂的喜悦创造了它们，因为祂是天主。作为万有的创造者、同议者和塑造者，祂生了圣言；当祂将这圣言怀在祂内，且对受造世界仍不可见时，祂使祂成为可见的；先发出声音，作为光明来自光明生了祂，祂将祂作为世界之主、祂自己的理智而揭示给世界；虽然祂先前只对祂自己可见，对受造世界不可见，但祂使祂成为可见，为叫世界在祂的显现中看见祂，并能得救。\par

11. 这样，在祂之外出现了另一位。但我说\Emph{另一位}，并非指有两位天主，而是如同光明来自光明，或水来自泉源，或光线来自太阳。因为只有一种能力，来自万有；圣父就是万有，这能力——圣言——就是从祂而来。这就是进入世界、并显现为天主子的理智。因此，万有都是藉着祂，只有祂是从圣父而来。那么，谁又一再引入众多神祇呢？因为所有人，尽管不情愿，都被迫承认这个事实：万有都归于一体。如果万有都归于一体，即使根据瓦伦廷、马西翁、塞林图和他们所有愚蠢的言论，他们也（尽管不情愿）被逼到这样的立场：他们必须承认那一位是万有的原因。这样，这些人虽然不愿意，也附和真理，承认一位天主按祂的善意创造了万有。祂赐下了法律和先知；在赐予时，祂使他们藉圣神发言，为叫他们蒙受圣父能力的灵感，得以宣告圣父的计划和旨意。\par

12. 圣言就这样在先知们中行动，谈论自己。因为祂早已成为自己的使者，并显示圣言将在人中间显现。为此，祂这样呼喊：\BibleQuote{我向不寻找我的人显现，被不寻求我的人找到。} 谁是所显现的那位，不就是圣父的圣言吗？——圣父派遣了祂，并在祂内将出自祂的能力显示给人。因此，圣言就这样显现了，正如有福的若望所说。因为他总结了先知们所说的话，表明这就是万有藉以受造的圣言。因为他这样说：\BibleQuote{在起初已有圣言，圣言与天主同在，圣言就是天主。万有是藉着祂而造成的，凡受造的，没有一样不是藉着祂造成的。} 下文又说：\BibleQuote{世界是藉着祂造成的，但世界不认识祂；祂来到自己的领域，自己的人却没有接受祂。} 所以他说，如果世界是藉着祂造成的，按照先知的话：\BibleQuote{诸天是藉着上主的圣言造成的，} 那么这就是那显现的圣言。因此，我们看见圣言成了血肉，藉着祂认识圣父，我们信圣子，钦崇圣神。让我们看看圣经关于圣言将来显现的宣告的见证。\par

13. 现在耶肋米亚说：\BibleQuote{谁曾站在上主的会议中，以看见並听祂的话？} 但天主的圣言是可见的，而人的言语是可听的。当他说看见圣言时，我必须相信这可见的圣言已被派遣。没有别的被派遣，只有圣言。伯多禄在向科尔乃略百夫长讲话时证实了祂被派遣：\BibleQuote{天主借耶稣基督——祂是万有的主——宣讲和平，将祂的圣言传给以色列子民。这就是万有之主的天主。} 因此，如果圣言是藉着耶稣基督派遣的，那么圣父的旨意就是耶稣基督。\par

14. 这些事，弟兄们，是由圣经所宣告的。蒙福的若望在他的 \BibleBook{若望}{福音} 的见证中，向我们叙述了这\OriginalTerm{安排}{economy (disposition)}，并承认这圣言是天主，他说：\BibleQuote{在起初已有圣言，圣言与天主同在，圣言就是天主。}那么，如果圣言与天主同在，并且也是天主，接着是什么呢？有人会说祂讲的是两个神吗？我确实不会说两个神，而是一个；但有两个位格，以及第三个\OriginalTerm{安排}{economy (disposition)}，即圣神的恩宠。因为圣父确实是一个，但有两个位格，因为还有圣子；然后还有第三位，圣神。圣父命令，圣言执行，圣子被显示，借着祂，圣父被人相信。和谐的安排归于一个天主；因为天主是独一的。命令的是圣父，听从的是圣子，赐予理解的是圣神：圣父是\Emph{超越万有}的，圣子是\Emph{贯通万有}的，圣神是\Emph{在万有内}的。我们无法以其他方式思想独一的天主，除非真实地相信父、子和圣神。因为犹太人光荣了（或夸耀）圣父，但未感谢祂，因为他们不承认圣子。门徒承认了圣子，但不是在圣神内；因此他们也否认了祂。因此，圣父的圣言知道这安排（处置）和圣父的旨意，即圣父只愿以此方式受敬拜，就在祂从死者中复活后给门徒这个命令：\BibleQuote{你们要去使万民成为门徒，因父、及子、及圣神之名给他们授洗。}借此祂表明，凡省略其中任何一位的，都没有完全光荣天主。因为正是通过这个天主圣三，圣父才得光荣。因为圣父愿意，圣子做了，圣神彰显了。整部圣经便宣告了这真理。\par

15. 但有人会对我说：当你称圣子为圣言时，你向我提出了一件陌生的事。因为若望确实讲到圣言，但那是一种修辞方式。\Emph{不，绝对不是修辞方式。}因为当他这样呈现这从起初就有、如今已被派遣的圣言时，他在 \BibleBook{}{默示录} 后面说：\BibleQuote{我看见天开了，看，有一匹白马；骑在祂上的（被称为）\InnerQuote{忠诚}和\InnerQuote{真实}，祂凭正义审判和作战。祂的眼睛如同火焰，祂头上戴着许多王冠；祂有一个名字写着，除了祂自己没有人知道。祂穿着浸过血的衣服；祂的名字被称为\InnerQuote{天主的圣言}。}那么，弟兄们，请看那蘸血的衣服如何以象征方式表示肉身，借着它，天主那\OriginalTerm{不受苦难的}{impassible}圣言遭受了苦难，正如先知们为我作证。因为蒙福的米该亚这样说：\BibleQuote{雅各伯家激怒了上主的神。这些是他们的追求。祂的话在他们中间不是好的吗，他们不是行得正吗？他们却起来敌对平安的圣容，剥去了祂的光荣。}那是指祂在肉身中的苦难。同样，蒙福的保禄也说：\BibleQuote{因为法律因肉性软弱而做不到的，天主却做到了：祂派遣了自己的儿子，带着罪恶肉身的形状，为了罪，在肉身上定了罪案，为使法律所要求的正义，成全在我们这些不随从肉身、而随从圣神的人身上。}那么，天主借着肉身派遣了自己的什么儿子，岂不就是圣言吗？祂称圣言为儿子，因为祂将来要成为（或被生）儿子。祂取了在人中表示亲情的普通名称而被称为儿子。因为圣言在降生成人之前、独自存在时，还不是完美的儿子，尽管祂是完美的圣言、独生子。肉身也不能离开圣言而自立存在，因为它的\OriginalTerm{自立存在}{subsistence}在圣言内。这样，一个完美的天主子便被彰显了。\par

16. 这些确实是关于圣言降生成人的见证；还有许多其他见证。但让我们也看看手头的问题——即，弟兄们，实际上，圣父的能力，即圣言，从天上下来，而不是圣父自己。因为祂如此说：\BibleQuote{我出自父，并且来到了。}那么，在这句话\InnerQuote{我出自父}中，主语除了圣言还有谁呢？从祂所生的又是什么，岂不就是圣神，即圣言吗？但你会对我说：祂是如何被生的？在你自身的事上，你无法给出你如何被生的解释，尽管你每天看到按人来说的原因；你也不能精确述说祂身上的\OriginalTerm{安排}{economy}。因为你没有能力去认识造物主那娴熟而不可言喻的技巧（方法），只能看到、理解并相信人是天主的工程。此外，你询问圣言受生的情形，而天主圣父按祂的善意随意生了祂。你得知天主创造了世界，还不够吗？你还敢问祂从何而造？你得知天主子为了你的救恩而显现给你，只要你相信，还不够吗？你还好奇地追问祂如何按圣神受生？事实上，只有两个人受委托述说祂按肉身的受生；而你竟敢寻求按圣神受生的叙述，那是圣父自己保留的，要在那时启示给圣人和那些配见祂面的人？你该满足于基督所说的这话：\BibleQuote{那由圣神生的就是神，}正如祂借着先知论到圣言的受生时，表明了祂受生的事实，但保留了方式和途径的问题，只在自己决定的时间启示。因为祂这样说：\BibleQuote{在晨星之前，我从母腹中生了祢。}\par

17. 这些见证对于信德而追求真理的人是足够的，不信的人不会相信任何见证。因为圣神确实藉宗徒的口为这一点作了见证，说：\Quote{谁相信了我们的报道呢？}因此，让我们不要表现出不信，免得这话应在我們身上。所以，亲爱的弟兄们，让我们按照宗徒的传统相信：天主圣言从天降下，（进入）童贞玛利亚，为的是从她取得肉躯，并取了一个人性，即理性的灵魂，从而成为完全的人，只是没有罪，好能拯救堕落的人，并赐予永生给那些相信祂名的人。因此，在所有事上，真理之言向我们显明，即：圣父是惟一的，祂的圣言与祂同在，万物由祂造成；正如我们前面所说，圣父在末世为拯救人类派遣了这圣言。这圣言由法律和先知预言将要来到世界。正如那时所预言的，祂也同样来到并显示自己，由童贞女和圣神成为新人；因为祂既有父的属天本性作为圣言，又有属地本性，经由童贞女从旧亚当取得肉躯，如今进入世界，在身体中显示为天主，也作为完美的人出现。因为祂成为人不是仅仅外表或转化，而是真实地。\par

18. 因此，虽然被证明为天主，祂也不拒绝属于人的条件，因为祂饥饿、劳苦、干渴、疲倦、因恐惧而逃跑、在患难中祈祷。祂作为天主有警醒的本性，却在枕头上睡觉。祂为此目的来到世界，却求免于受苦的杯。在痛苦中祂汗流如血，并被一位天使加添力量，而祂自己却坚固那些相信祂的人，并以祂的工作教导人蔑视死亡。祂知道犹大是何种人，却被犹大出卖。祂先前被犹大尊为天主，现在却被盖法藐视。祂被黑落德蔑视，而祂自己将要审判全地。祂被比拉多鞭打，祂自己担当了我们的病弱。祂被士兵戏弄，而祂的命令下有千千万万的天使和总领天使侍立。祂将穹苍固定如穹顶，却被犹太人钉在十字架上。祂与圣父不可分离，却向圣父呼喊，并把祂的灵魂交托给祂；低头，祂交付了灵魂，祂曾说过：\Quote{我有权舍弃我的生命，也有权再取回它；}因为死亡不能胜过祂，祂本身就是生命，祂这样说：\Quote{我自己舍弃它。}祂慷慨地将生命赐给众人，肋旁却被长矛刺透。祂复活死人，却被裹在殓布里放在坟墓中，第三日由圣父复活，而祂自己就是复活和生命。因为这一切祂为我们完成了，祂为了我们成了和我们一样的人。因为\Quote{祂亲自担当了我们的病弱，背负了我们的疾病；为了我们的缘故祂受折磨，}正如先知依撒意亚所说。这就是天使歌颂、牧羊人看见、西默盎期待、亚纳证明的那一位。这就是贤士们寻求、星星指示的那一位；祂在父的殿中，由若翰指出，并由圣父从上用声音作证：\Quote{这是我的爱子，你们要听从祂。}祂战胜了魔鬼，戴着胜利的冠冕。这就是纳匝肋人耶稣，祂被邀请到加纳的婚宴，将水变成酒，斥责因狂风翻腾的海，在海面上步行如履平地，使生来盲人看见，叫死了四天的拉匝禄复活，行了许多大能的事，赦免罪过，赐能力给门徒，被长矛刺透时从神圣的肋旁流出血和水。为了祂，太阳变暗，白昼无光，磐石崩裂，帐幔撕裂，大地的根基摇动，坟墓打开，死人复活，统治者羞愧，因为他们看见宇宙的统御者在十字架上闭目交魂。受造物看见，就惊慌；因不能承受祂极度荣耀的景象，就用黑暗遮掩自己。祂向门徒嘘气，赐给他们圣神，在门关着时来到他们中间，在门徒注视中被云彩接到天上，坐在圣父的右边，将来要作为生者死者的审判者再来。这就是为了我们成为人的天主，圣父也将万有服在祂脚下。愿光荣与权能归于祂，和圣父及圣神，在圣教会中，从今世直到永远。阿们。\par

\SourceRecord{Translated by J.H. MacMahon.；Ante-Nicene Fathers；Vol. 5.；Edited by Alexander Roberts, James Donaldson, and A. Cleveland Coxe.；Buffalo, NY: Christian Literature Publishing Co.,；1886.；<http://www.newadvent.org/fathers/0521.htm>.}
